\section{Introduction}
\label{sec:intro}

Ptolemy's \emph{Almagest} \cite{Toomer1998} predicted planetary positions
well enough to serve as the working astronomy of the Mediterranean
and Islamic worlds for well over a thousand years, using circles carried by other circles
\cite{Neugebauer1975,Evans1998}, with motion that is uniform about the circle's
centre or, for the equant, about an offset point (Sec.~\ref{sec:history}). That a scheme so alien to modern dynamics
could be so accurate invites a modern explanation. A popular one is that
epicycles are ``just Fourier series'': a sum of rotating
vectors can trace any closed curve, so of course it fits the sky
\cite{Hanson1960,Mathologer2018,3b1b2019,WikiDeferent}. The statement is
true in spirit but, as usually made, it is loose in three ways. First, the
motion of a planet seen from Earth involves at least two generally
incommensurate frequencies, the mean motions of the planet and of the Earth,
so the geocentric path is not periodic and has no Fourier series in the
strict sense; the appropriate setting is that of Bohr's almost-periodic
functions \cite{Bohr1925}. Second, Ptolemy's actual models use eccentrics
and the equant and are not harmonic expansions
\cite{Neugebauer1975,Evans1984}. Third, if epicycles can fit anything, the
representation is empirically empty \cite{Carman2010}, unless the radii and
rates are fixed by something other than the data.

In this paper we ask a precise version of the question. Given a body on a
Newtonian orbit, what is the epicycle system, with radii and angular rates
determined by the orbital elements rather than fitted, that reproduces its
geocentric position, how many circles does it need for a prescribed accuracy,
and for which motions does the answer remain valid? The answer has two
parts. For unperturbed Kepler motion the heliocentric position is an exact
Fourier series in the mean anomaly with Bessel-function coefficients
\cite{Dutka1995,Watson1944,MurrayDermott1999}. Subtracting the Earth's
position gives a sum of two such series, one for the body and one
for the Earth, which is an epicycle system in which every radius and rate is
a function of $(a,e,i,\omega,\Omega,M_0)$ of the two orbits. For perturbed
motion the representation is exact only on analytic invariant tori (whose
existence for a given trajectory is not guaranteed); for other trajectories it is
an empirical finite-window approximation, whose frequencies and amplitudes can be
recovered numerically by Laskar's frequency analysis
\cite{Laskar1990,Laskar1992,Laskar1999}.

\paragraph*{What is, and is not, claimed.}
The correspondence between epicycles and Fourier or quasi-periodic expansions
is not new, and we do not claim to have discovered it. It is stated
explicitly by Gallavotti \cite{Gallavotti2001}; the Bessel expansion of
Kepler motion goes back to Lagrange and Bessel
\cite{Dutka1995,Kapteyn1893}; and modern planetary theories are, in
practice, enormous sums of uniform circular motions
\cite{Bretagnon1988,Kudryavtsev2016}. Section~\ref{sec:prior} surveys this
literature. What we found no source to do together, after a search described
there, is the following, and these are the contributions of this paper.
\begin{enumerate}
\item An explicit, closed-form, three-dimensional geocentric epicycle
expansion of an arbitrary Keplerian body, with Bessel (Kapteyn) amplitudes,
built from the decomposition of each ellipse into two counter-rotating
circles, with inclination handled exactly as uniformly rotating vectors in
tilted planes.
\item Truncation-error bounds for this expansion, and the resulting number
of epicycles needed for a given accuracy as a function of eccentricity, for
both the body and the Earth terms.
\item A numerical frequency-analysis (NAFF) study. On unperturbed Keplerian
geocentric trajectories (Mars, Eros), NAFF recovers the analytic Bessel
epicycle terms (for the ten leading Mars terms over 200 years, with 10-day sampling and 24 lines, relative frequency error
$2\times10^{-10}$ and relative radius error $9\times10^{-8}$).
On an N-body-perturbed Ceres trajectory (Sun, Earth, Jupiter and
Saturn; 200 years), the NAFF reconstruction error is compared with that of
least-squares-fitted six-element Kepler orbits, fitted to the same heliocentric
data with matched weights and with the same integrated Earth subtracted from every
reconstruction; NAFF is clearly better than the fitted Kepler orbit only for
$K\gtrsim50$ lines and only in the interior of the window, where edge effects
dominate the full window.
\item A precise statement of validity: the representation is exact for Kepler motion and
on analytic invariant tori, an empirical finite-window approximation for other
perturbed motion (including chaotic motion), and does not apply to unbound
orbits.
\end{enumerate}
In addition we provide an open-source interactive three-dimensional
visualization of the resulting epicycle chains as a pedagogical tool. The
claim of novelty is about our search, not a proof of absence.

\paragraph*{Why this is not curve fitting.}
In the construction below, every radius and angular rate is determined by the
orbital elements of the body and of the Earth. Epicycles thus become a
derived coordinate representation of Newtonian motion rather than a fitting
model, which is our answer to the objection that an epicycle system
explains nothing because it can reproduce anything.

\paragraph*{Outline.}
Section~\ref{sec:history} gives brief historical background and
Sec.~\ref{sec:prior} reviews prior work and positions our contribution.
Section~\ref{sec:theory} derives the expansion and its truncation bounds,
and Sec.~\ref{sec:numerics} describes the numerical methods (Kepler
propagation, N-body integration, and NAFF). Section~\ref{sec:results} reports
the results, Sec.~\ref{sec:viz} describes the visualization, and
Sec.~\ref{sec:discussion} discusses the domain of validity before
Sec.~\ref{sec:conclusion} concludes. Bessel-function identities used in the
derivation are collected in the Appendix.
